\section{Evaluation}
\label{sec:evaluation}

We evaluate \system{} by answering four questions.  First, can heterogeneous
execution sustain the 30-FPS real-time target once communication and runtime
overheads are included?  Second, does DAQ preserve metric-depth quality on
challenging sensor data?  Third, how does the complete system compare with
recent depth-completion methods?  Finally, what FPGA resource, communication,
power, and energy costs are required to obtain these gains?

\subsection{Methodology}

\paragraph{Baselines.}
We use two groups of baselines for different purposes.  The controlled system
comparison runs the same depth-completion network in three deployments:
\emph{GPU-only} on an RTX~4090, \emph{Thor-only} on the target edge SoC, and
\system{} with its encoder on Thor and decoder on FPGA.  The GPU-only and
Thor-only baselines execute the complete model sequentially on one device;
\system{} includes feature transfer and output reconstruction in its
end-to-end time.  For the component study, \system{} without DAQ retains the
same partition, asynchronous runtime, accelerator, and target precisions, but
replaces the proposed depth-aware policy with conventional low-precision
quantization.  This makes the comparisons in Table~\ref{tab:methodology}
attributable to deployment and quantization choices rather than to different
networks.

For model-level context, we additionally compare against CDM, PriorDA,
PromptDA, and OMNI-DC, following the methods introduced in
Section~\ref{sec:introduction}.  We use their released checkpoints and official
evaluation procedures whenever available.  Because their published FPS values
come from different GPU platforms, the state-of-the-art throughput plot labels
the execution platform of every method and does not treat cross-platform FPS
as a controlled speedup claim.

\paragraph{Benchmarks.}
We evaluate all accuracy-oriented experiments on HAMMER and ClearPose.  These
datasets exercise different object geometries and difficult sensing conditions,
so results are reported for each dataset rather than averaged into one score.
Every method uses the same test split, validity mask, crop, input resolution,
and depth range within a dataset.  When a released method requires a different
input representation, only the documented preprocessing is changed; the
ground-truth mask and metric computation remain common.  The system and
ablation experiments use batch size one, which matches continuous RGB-D
capture on the target robotic platform.

\paragraph{Experiment platform.}
The prototype comprises an NVIDIA Jetson AGX Thor and an Intel Stratix FPGA
connected through the system interface described in Section~\ref{sec:framework}.
The FPGA clock, external-memory configuration, link rate, Thor power mode, and
software versions are held fixed across the two \system{} variants.  The
RTX~4090 baseline uses the same model weights, input resolution, preprocessing,
and output validation.  Before submission, the final platform table will state
the exact CUDA/TensorRT stack, FPGA device and achieved clock, host--FPGA link,
and precision assigned to every model stage.

\paragraph{Metrics and measurement protocol.}
We measure both single-frame latency and steady-state throughput.  Latency is
decomposed into preprocessing, encoder execution, boundary-tensor transfer,
FPGA decoding, return transfer, and output reconstruction; FPS is measured only
after the asynchronous pipeline is filled.  Timing includes synchronization
needed for a valid completed-depth output, rather than reporting kernel time
alone.  Each result will state the warm-up length, measured frame count, and
run-to-run variation.

Depth quality is evaluated with AbsRel, RMSE, MAE, and $\delta_{1.05}$.  Lower
is better for the three error metrics and higher is better for
$\delta_{1.05}$.  We report system power, energy per completed frame, and FPS/W
under a consistent measurement boundary.  FPGA resource results use post-fit
utilization and achieved frequency rather than pre-synthesis estimates.

\begin{table}[t]
  \caption{Configurations used for the controlled system comparison and
  component ablation.  ``Conventional LP'' and ``DAQ'' use the same target
  bit widths; they differ only in the quantization policy.}
  \label{tab:methodology}
  \centering
  \small
  \setlength{\tabcolsep}{3.0pt}
  \renewcommand{\arraystretch}{1.10}
  \begin{tabular}{lccc}
    \toprule
    Configuration & Encoder & Decoder & Precision policy \\
    \midrule
    RTX 4090 only & RTX 4090 & RTX 4090 & reference \\
    Thor only & Thor & Thor & reference \\
    \system{} w/o DAQ & Thor & FPGA & conventional LP \\
    \system{} w/ DAQ & Thor & FPGA & depth-aware LP \\
    \bottomrule
  \end{tabular}
\end{table}

\subsection{Ablation Experiments}

Figure~\ref{fig:ablation-results} evaluates the four configurations in
Table~\ref{tab:methodology}.  We interpret the bars cumulatively: the
single-device baselines establish the deployment cost, the system without DAQ
exposes the benefit of heterogeneous execution, and the complete system shows
whether depth-aware precision recovers quality without giving up the hardware
path.

\paragraph{Throughput and heterogeneous execution.}
Figure~\ref{fig:ablation-results}(a) reports absolute FPS and marks the 30-FPS
real-time threshold.  Thor-only versus \system{} without DAQ is the primary
controlled comparison for the Thor--FPGA partition: both process the same
frames with the same network, while only the latter offloads the decoder and
overlaps encoder, communication, and decoder work across frames.  We therefore
report not only their FPS difference, but also the stage-latency breakdown and
the fraction of transfer hidden by the asynchronous schedule.  \system{} with
and without DAQ is then compared under the same runtime to determine whether
the precision policy changes throughput.  GPU-only is retained as a high-end
performance reference, not as an edge-device energy baseline.  The completed
discussion will state absolute FPS first, followed by speedup over Thor-only,
and will explain any remaining gap using the slowest pipeline stage rather
than attributing it generically to ``offloading.''

\paragraph{Metric-depth accuracy.}
Figures~\ref{fig:ablation-results}(b)--(e) separate HAMMER and ClearPose and
show both an error metric and a threshold accuracy metric.  The decisive
comparison is \system{} without versus with DAQ because their model partition,
target bit widths, accelerator, and runtime are identical.  Any quality
difference therefore measures the value of coordinating the quantized encoder
MLP, boundary representation, and FPGA decoder using metric-depth supervision.
The GPU-only and Thor-only results provide the unpartitioned reference and
reveal the total degradation introduced by the heterogeneous low-precision
path.  We report each dataset independently and discuss whether improvements
come from reducing large depth errors, increasing the fraction of accurate
pixels, or both; a single aggregate ``near-SOTA'' statement is not used in
place of these measurements.

% The component ablation is the only full-width figure in the paper.
\begin{figure*}[t]
  \centering
  \figslot{figures/ablation-results.pdf}{1.62in}{%
    Results of ablation experiments --- grouped bar charts\\[5pt]
    AS-Depth w/ DAQ / AS-Depth w/o DAQ / GPU only / Thor only\\
    (a) FPS $\uparrow$ with 30-FPS line \quad
    (b) HAMMER AbsRel $\downarrow$ \quad
    (c) HAMMER $\delta_{1.05}$ $\uparrow$\\
    (d) ClearPose AbsRel $\downarrow$ \quad
    (e) ClearPose $\delta_{1.05}$ $\uparrow$}
  \caption{Results of the component ablation.  Throughput isolates the effect
  of heterogeneous execution, while the HAMMER and ClearPose panels quantify
  the accuracy recovered by DAQ under the same partition and target precision.}
  \Description{Five grouped bar charts comparing FPS, HAMMER absolute
  relative error and delta accuracy, and ClearPose absolute relative error and
  delta accuracy for AS-Depth with DAQ, AS-Depth without DAQ, GPU-only, and
  Thor-only.}
  \label{fig:ablation-results}
\end{figure*}

\subsection{Comparison with State-of-the-Art Methods}

\paragraph{Throughput.}
Figure~\ref{fig:sota-results}(a) compares the reported inference rate of CDM,
PriorDA, PromptDA, OMNI-DC, and \system{}.  Each bar is annotated with its
execution platform and the plot includes the 30-FPS reference line.  This
comparison answers a deployment question--whether current depth-completion
methods can service a real-time sensor stream--while the controlled
GPU-only/Thor-only comparison above supports architectural speedup claims.  In
the completed result discussion, we report \system{}'s measured end-to-end FPS
and identify whether its advantage comes from the coarse model partition,
cross-frame overlap, or low-precision execution, using the ablation evidence
rather than the cross-platform bars alone.

\paragraph{Depth accuracy.}
Figures~\ref{fig:sota-results}(b) and (c) compare all methods under common
HAMMER and ClearPose protocols.  The discussion reports the absolute metric
for every method, \system{}'s gap to the best-performing method, and its rank on
each dataset.  We then relate that gap to the deployment result: the relevant
question is whether \system{} reaches real-time edge throughput while keeping
the accuracy loss small and consistent across both datasets.  Claims of
state-of-the-art or near-state-of-the-art quality will be made only for metrics
and datasets for which the measured values support them.

\begin{figure}[t]
  \centering
  \figslot{figures/sota-results.pdf}{1.68in}{%
    Results of comparison with state of the art --- grouped bar charts\\[5pt]
    (a) FPS $\uparrow$: CDM / PriorDA / PromptDA / OMNI-DC / AS-Depth,\\
    with platform labels and 30-FPS line\\
    (b) HAMMER accuracy \qquad (c) ClearPose accuracy}
  \caption{Comparison with state-of-the-art depth-completion methods.
  Throughput is labeled by platform and provides deployment context; accuracy
  is evaluated on HAMMER and ClearPose with a common protocol.}
  \Description{Three grouped bar charts comparing FPS, HAMMER depth accuracy,
  and ClearPose depth accuracy across CDM, PriorDA, PromptDA, OMNI-DC, and
  AS-Depth.}
  \label{fig:sota-results}
\end{figure}

\subsection{Overhead Analysis}

\paragraph{FPGA resource and communication overhead.}
Table~\ref{tab:hardware-overhead} reports post-fit logic, register, DSP, and
on-chip-memory utilization together with achieved frequency.  To connect
resource cost to the system design, we also report boundary bytes per frame,
measured link latency, and standalone decoder latency.  These measurements
show whether the selected precision fits the target device and whether feature
transfer is sufficiently small to be hidden behind Thor encoding.  Resource
percentages are always reported with the device capacity and synthesis tool
settings so that raw counts are interpretable.

\paragraph{Power and energy efficiency.}
Power is measured for Thor-only, the active FPGA, and the complete
Thor--FPGA system using the same input stream and steady-state window used for
FPS.  We derive energy per frame as system power divided by completed frames
per second and additionally report FPS/W.  The analysis separates the extra
power introduced by the FPGA and link from the energy saved by shortening the
end-to-end execution time.  RTX~4090 power is reported as a separate reference
because its platform boundary differs from the embedded prototype.  The final
discussion will pair every energy-efficiency ratio with both its power and
throughput components, explaining when additional hardware power is offset by
a shorter end-to-end execution time.

\begin{table}[t]
  \caption{Post-fit hardware, communication, and energy overhead.  Values
  marked ``--'' await implementation or measurement.}
  \label{tab:hardware-overhead}
  \centering
  \small
  \setlength{\tabcolsep}{3.0pt}
  \renewcommand{\arraystretch}{1.08}
  \begin{tabular}{lr@{\qquad}lr}
    \toprule
    Metric & Value & Metric & Value \\
    \midrule
    Logic / registers & -- & DSP blocks & -- \\
    On-chip memory & -- & Frequency & -- \\
    Feature bytes/frame & -- & Link latency & -- \\
    Decoder latency & -- & System power & -- \\
    Energy/frame & -- & FPS/W & -- \\
    \bottomrule
  \end{tabular}
\end{table}
